% chapters/03/section-03.tex
\section{Memasang Pygame dan Memeriksa Pemasangan}
\label{sec:memasang-pygame}

\subsection{Kedudukan Pygame}

Pygame bukan bahasa pemrograman dan bukan pula program yang berdiri sendiri.
Ia adalah \textbf{pustaka}, yaitu kumpulan fungsi yang dapat Anda panggil dari
program Python untuk mengurus pekerjaan yang berkaitan dengan permainan.
Pekerjaan itu meliputi membuka jendela, menggambar bentuk dan teks, memuat
berkas gambar dan bunyi, serta membaca masukan pemain \cite{pygame-docs}.

Pekerjaan tingkat rendah di bawah Pygame ditangani oleh pustaka lain bernama
SDL (\istilah{Simple DirectMedia Layer}). Program Anda tidak memanggil SDL
secara langsung; Pygame yang menjadi perantaranya. Susunan lapisan tersebut
diperlihatkan pada Gambar~\ref{fig:lapisan-perkakas}.

\begin{figure}[htbp]
\centering
% chapters/03/figures/lapisan-perkakas.tikz
% Lapisan perkakas: program Anda memanggil Pygame, Pygame memakai SDL,
% SDL berbicara dengan sistem operasi dan perangkat.
\begin{tikzpicture}[
  lapis/.style={
    rectangle, rounded corners=2pt, draw=ubPrimary, thick, fill=ubAlt,
    align=center, text width=70mm, inner sep=4pt, minimum height=10mm,
    font=\footnotesize
  },
  lapissdl/.style={lapis, draw=ubAccent, fill=white},
  lapisos/.style={lapis, draw=ubLine, fill=ubGray, text=black},
  panah/.style={-{Stealth[length=2.2mm]}, thick, draw=ubAccent},
  sisi/.style={font=\scriptsize\itshape, ubPrimary, align=left}
]

\node[lapis] (program) at (0,0)
  {\textbf{Program Anda}\\[1pt]\scriptsize kode yang Anda tulis pada folder \code{code/}};
\node[lapissdl] (pygame) at (0,-1.8)
  {\textbf{Pygame}\\[1pt]\scriptsize jendela, gambar, huruf, bunyi, dan penanganan peristiwa};
\node[lapissdl] (sdl) at (0,-3.6)
  {\textbf{SDL}\\[1pt]\scriptsize lapisan penerjemah ke perangkat keras};
\node[lapisos] (os) at (0,-5.4)
  {\textbf{Sistem operasi dan perangkat}\\[1pt]\scriptsize layar, papan tombol, tetikus, pengeras suara};

\draw[panah] (program) -- (pygame);
\draw[panah] (pygame) -- (sdl);
\draw[panah] (sdl) -- (os);

\node[sisi, right=5mm of program, text width=42mm, anchor=west]
  {perintah yang Anda tulis: \code{pygame.display.set_mode}};
\node[sisi, right=5mm of sdl, text width=42mm, anchor=west]
  {diterjemahkan ke perintah perangkat};
\end{tikzpicture}

\caption{Lapisan perkakas dari kode Anda sampai perangkat keras}
\label{fig:lapisan-perkakas}
\end{figure}

Susunan berlapis itu menjelaskan dua hal yang akan Anda jumpai nanti. Pertama,
pesan kesalahan kadang menyebut SDL, bukan Pygame, karena masalahnya berada di
lapisan yang lebih dalam. Kedua, versi SDL yang terpasang ikut menentukan
kemampuan yang tersedia, sehingga nomor versinya perlu Anda catat ketika
meminta bantuan.

\subsection{Memasang Pygame}

Pemasangan dilakukan dengan \istilah{pip}, yaitu pengelola paket milik Python.
Jalankan perintah berikut pada terminal.

\begin{kodekeluaran}
python -m pip install pygame
\end{kodekeluaran}

Perintah tersebut dibaca sebagai berikut: ``jalankan modul \code{pip} dengan
penafsir \code{python} yang sedang aktif, lalu pasang paket bernama
\code{pygame}''. Bentuk \code{python -m pip} sengaja dipakai alih-alih
\code{pip} saja. Alasannya, apabila pada komputer terdapat lebih dari satu
versi Python, perintah \code{pip} dapat menunjuk versi yang berbeda dengan
\code{python}. Akibatnya Pygame terpasang pada penafsir yang tidak Anda pakai,
dan program tetap melaporkan bahwa Pygame tidak ditemukan. Bentuk
\code{python -m pip} memastikan pemasangan dilakukan pada penafsir yang sama
dengan yang akan menjalankan program Anda.

Apabila pemasangan gagal karena masalah izin, tambahkan pilihan
\code{--user} sehingga paket dipasang di folder pribadi Anda, bukan di
folder sistem.

\begin{kodekeluaran}
python -m pip install --user pygame
\end{kodekeluaran}

\subsection{Memeriksa Pemasangan dengan Skrip}

Menjalankan satu perintah pemasangan belum membuktikan bahwa lingkungan Anda
siap. Perlu dipastikan pula bahwa Pygame dapat dipanggil, pengandar video
tersedia, dan huruf bawaan dapat dimuat. Untuk keperluan itu, buku ini
menyediakan skrip pemeriksaan yang dapat Anda pakai berulang kali, misalnya
ketika berpindah komputer.

\kodepython{python/ch03_periksa_lingkungan.py}{Skrip pemeriksaan kesiapan lingkungan kerja}{lst:ch03-periksa}

Jalankan skrip tersebut dari terminal, dan Anda akan memperoleh keluaran
serupa berikut.

\begin{kodekeluaran}
============================================================
 PEMERIKSAAN LINGKUNGAN KERJA
 Pengantar Pemrograman Game
============================================================

Python           : 3.11.9 (siap)
Sistem operasi   : Windows 10 AMD64
Penafsir Python  : C:\Program Files\Python311\python.exe
Pygame           : 2.6.1
SDL              : 2.28.4
Pengandar video  : windows
Huruf bawaan     : tersedia (tinggi 16 piksel)

Hasil: lingkungan siap dipakai. Lanjutkan ke praktikum.
\end{kodekeluaran}

Baris \istilah{Penafsir Python} menunjukkan berkas penafsir yang sedang dipakai.
Catatlah baris tersebut, karena ketika muncul masalah ``Pygame tidak
ditemukan'', baris itulah yang menunjukkan penafsir mana yang harus
diperiksa. Baris \istilah{Pengandar video} menunjukkan cara Pygame berbicara
dengan layar; pada komputer biasa nilainya \code{windows} atau \code{x11},
sedangkan pada pengujian tanpa layar nilainya \code{dummy}.

Skrip tersebut juga memberikan kode keluar yang dapat diperiksa oleh perkakas
lain: bernilai 0 apabila lingkungan siap, dan 1 apabila ada yang belum siap.
Cara itu dipakai agar pemeriksaan dapat dijalankan berulang kali tanpa harus
membaca seluruh keluarannya.

\begin{catatabox}
Saat pertama kali dijalankan, Pygame menampilkan dua baris sambutan pada
terminal, yaitu keterangan versi dan ajakan mengunjungi laman komunitas
Pygame. Kedua baris itu bukan kesalahan dan bukan bagian dari keluaran
program Anda. Selain itu, pada sebagian versi muncul peringatan mengenai
paket \code{pkg_resources} yang sudah usang. Peringatan tersebut berasal dari
Pygame sendiri, tidak memengaruhi jalannya program, dan boleh diabaikan.
Perbedaan antara peringatan dan kesalahan dijelaskan pada
Bagian~\ref{sec:menelusuri-kesalahan}.
\end{catatabox}

\subsection{Mengatasi Masalah Pemasangan}

Masalah pada tahap ini jarang berkaitan dengan kode, dan hampir selalu
berkaitan dengan penafsir yang dipakai. Tabel~\ref{tab:masalah-pemasangan}
memuat gejala yang paling sering muncul beserta penyebab dan tindakannya.

\begin{table}[htbp]
\centering
\caption{Masalah pemasangan yang sering terjadi dan penanganannya}
\label{tab:masalah-pemasangan}
{\setlength{\tabcolsep}{4pt}\footnotesize%
\begin{tabularx}{\textwidth}{@{}L{4.0cm} L{4.0cm} Y@{}}
\toprule
\headrow \textbf{Gejala} & \textbf{Penyebab yang paling sering} & \textbf{Tindakan} \\
\midrule
Perintah \code{python} tidak dikenali &
Python belum terpasang, atau tidak terdaftar pada \istilah{PATH} sistem. &
Pasang ulang dengan pilihan ``Add python.exe to PATH'' dicentang, atau pakai
perintah \code{py} sebagai gantinya. \\
\addlinespace
Perintah \code{pip} tidak dikenali &
\code{pip} tidak terdaftar pada \istilah{PATH}, walaupun Python terpasang. &
Selalu pakai bentuk \code{python -m pip}, yang tidak bergantung pada
\istilah{PATH}. \\
\addlinespace
\code{ModuleNotFoundError: No module named 'pygame'} &
Pygame belum terpasang, atau terpasang pada penafsir Python yang berbeda. &
Jalankan skrip pemeriksaan, perhatikan baris \istilah{Penafsir Python}, lalu
pasang ulang dengan \code{python -m pip install pygame}. \\
\addlinespace
\code{Access is denied} ketika memasang &
Tidak ada izin menulis pada folder instalasi Python. &
Tambahkan pilihan \code{--user} pada perintah pemasangan. \\
\addlinespace
\code{pygame.error: No available video device} &
Program dijalankan pada lingkungan tanpa layar, misalnya sesi jarak jauh. &
Kerjakan pada komputer laboratorium; untuk pengujian tanpa layar, setel
variabel \code{SDL_VIDEODRIVER} bernilai \code{dummy}. \\
\addlinespace
Jendela terbuka lalu langsung tertutup &
Program dijalankan dengan pilihan \code{--auto-close} yang mengirim sendiri
peristiwa penutup jendela. &
Periksa argumen yang dipakai pada baris perintah dan jalankan tanpa pilihan
tersebut. \\
\bottomrule
\end{tabularx}}
\end{table}
